\documentclass[11pt,a4paper]{article}
\usepackage[spanish,es-noquoting]{babel}
\usepackage[T1]{fontenc}
\usepackage[utf8]{inputenc}
\usepackage{lmodern}
\usepackage{microtype}
\usepackage[a4paper,margin=2.35cm,headheight=14pt]{geometry}
\usepackage{amsmath,amssymb}
\usepackage{booktabs,tabularx,array}
\usepackage{siunitx}
\sisetup{output-decimal-marker={,},per-mode=symbol}
\usepackage{tikz}
\usepackage{circuitikz}
\usepackage{enumitem}
\usepackage{fancyhdr}
\usepackage{listings}
\usepackage[hidelinks]{hyperref}
\setlength{\parindent}{0pt}
\setlength{\parskip}{0.42em}
\setlist[itemize]{leftmargin=1.5em,itemsep=0.12em,topsep=0.15em}
\setlist[enumerate]{leftmargin=1.5em,itemsep=0.12em,topsep=0.15em}
\renewcommand{\arraystretch}{1.16}
\pagestyle{fancy}
\fancyhf{}
\fancyfoot[C]{\thepage}
\renewcommand{\headrulewidth}{0pt}
\lstset{basicstyle=\ttfamily\small,frame=single,framerule=0.5pt,columns=fullflexible,keepspaces=true,showstringspaces=false,breaklines=true}
\newcommand{\encabezado}[2]{%
  {\Large\bfseries #1\par}
  \vspace{0.25em}
  {\normalsize #2\par}
  \vspace{0.5em}\hrule\vspace{0.6em}
}
\begin{document}
\encabezado{Detector de paso por infrarrojos}{Validación inicial del receptor y de la medida diferencial -- 19/11/2026}

\begin{tabularx}{\linewidth}{@{}lX@{}}
\textbf{Proyecto:} & \texttt{ir-doorway-detector}\\
\textbf{Placa:} & ESP32-S3 N16R8 (perfil de desarrollo DevKitC-1)\\
\textbf{Conexión:} & USB CDC; monitor serie a 115200 baudios\\
\textbf{Estado:} & Ensayo óptico en protoboard. Aún no se ha probado la detección en la puerta ni la PCB final.\\
\end{tabularx}

\section{Objetivo de la prueba}
En esta sesión quise comprobar si el sistema es capaz de distinguir el haz libre de un objeto que interrumpe la luz infrarroja. Repetí la prueba con una lámpara, con la iluminación normal de casa y con las luces apagadas. El objetivo era ver cuánto cambia la lectura al bloquear el haz y si la iluminación afecta a esa separación.

El montaje utiliza el LED infrarrojo L-53F3C controlado desde el GPIO4 mediante un BC337. El fototransistor BPW40 se conecta a la entrada analógica GPIO5 con una resistencia pull-up de \SI{10}{\kilo\ohm}. El circuito completo y los cálculos de polarización están documentados en el informe de circuitos y cálculos.

\section{Método de medida}
El programa realiza dos medidas consecutivas del ADC: una con el LED IR apagado y otra con el LED encendido. Para cada estado calcula la media de diez lecturas. Después obtiene la diferencia:
\begin{equation}
\overline{A}=\frac{1}{10}\sum_{i=1}^{10}A_i,\qquad
\Delta A=\overline{A}_{\mathrm{sin\ IR}}-\overline{A}_{\mathrm{con\ IR}}.
\label{eq:delta}
\end{equation}

El ADC está configurado a 12 bits, por lo que devuelve valores de 0 a 4095. El circuito receptor tiene una resistencia hacia 3,3 V: cuando el BPW40 recibe más luz, conduce más y la tensión en GPIO5 disminuye. Por eso, si el haz llega al fototransistor, espero obtener una diferencia $\Delta A$ elevada. Cuando coloco un objeto entre ambos, esa diferencia debe reducirse.

\begin{center}
\begin{tikzpicture}[>=latex, every node/.style={font=\small}]
\node[draw,minimum width=2.8cm,minimum height=0.9cm,align=center] (tx) at (0,0) {Emisor IR\\GPIO4};
\node[draw,minimum width=2.8cm,minimum height=0.9cm,align=center] (rx) at (6.8,0) {Receptor BPW40\\GPIO5 (ADC)};
\draw[->,dashed] (tx.east) -- node[above]{Haz infrarrojo} (rx.west);
\node[draw,minimum width=3.1cm,minimum height=0.83cm,align=center] (mcu) at (3.4,-1.65) {ESP32-S3\\$\Delta A=A_{\rm off}-A_{\rm on}$};
\draw[->] (mcu.west) -| (tx.south);
\draw[->] (rx.south) |- (mcu.east);
\end{tikzpicture}
\end{center}

\section{Resultados}
La tabla recoge los intervalos observados de $\Delta A$, tanto con el haz libre como con un objeto interpuesto. Son intervalos de lectura anotados en el ensayo; no representan un estudio estadístico ni incluyen un número de repeticiones documentado.

\begin{center}
\begin{tabularx}{\linewidth}{@{}Xrrr@{}}
\toprule
\textbf{Iluminación} & \textbf{Haz libre} & \textbf{Bloqueado} & \textbf{Separación mínima}\\
\midrule
Lámpara & 3378--3385 & 323--638 & 2740\\
Luz ambiental de casa & 3917--3924 & 391--604 & 3313\\
Oscuridad & 3913--3916 & 319--739 & 3174\\
\bottomrule
\end{tabularx}
\end{center}

En cada condición he calculado la separación tomando el menor valor con el haz libre y el mayor con el haz bloqueado:
\begin{equation}
M=\min(\Delta A_{\mathrm{libre}})-\max(\Delta A_{\mathrm{bloqueado}}).
\label{eq:margen}
\end{equation}
Por ejemplo, en la prueba con lámpara:
\begin{equation}
M_{\mathrm{lampara}}=3378-638=2740\ \text{cuentas ADC}.
\end{equation}
Si se comparan todos los ensayos a la vez, sin restringirse a una iluminación, la separación más conservadora es:
\begin{equation}
M_{\mathrm{global}}=3378-739=2639\ \text{cuentas ADC}.
\end{equation}
No hay solapamiento entre los dos estados dentro de los intervalos registrados.

\begin{figure}[htbp]
\centering
\begin{tikzpicture}[x=0.0025cm,y=0.60cm, every node/.style={font=\scriptsize}]
\foreach \x in {0,1000,2000,3000,4000} {
    \draw[dotted] (\x,0.45) -- (\x,6.5);
    \node[anchor=north] at (\x,0.42) {\x};
}
\foreach \y/\label in {6/{Lámpara: libre},5/{Lámpara: bloqueado},4/{Ambiente: libre},3/{Ambiente: bloqueado},2/{Oscuridad: libre},1/{Oscuridad: bloqueado}} {
    \node[anchor=east] at (-130,\y) {\label};
}
\draw[very thick] (3378,6) -- (3385,6); \fill (3381.5,6) circle (2.6pt);
\draw[very thick] (323,5) -- (638,5); \draw (323,4.87) -- (323,5.13); \draw (638,4.87) -- (638,5.13);
\draw[very thick] (3917,4) -- (3924,4); \fill (3920.5,4) circle (2.6pt);
\draw[very thick] (391,3) -- (604,3); \draw (391,2.87) -- (391,3.13); \draw (604,2.87) -- (604,3.13);
\draw[very thick] (3913,2) -- (3916,2); \fill (3914.5,2) circle (2.6pt);
\draw[very thick] (319,1) -- (739,1); \draw (319,0.87) -- (319,1.13); \draw (739,0.87) -- (739,1.13);
\node[anchor=north] at (2050,-0.1) {Diferencia ADC $\Delta A$};
\end{tikzpicture}
\caption{Intervalos observados en las tres condiciones. Los puntos representan los intervalos muy estrechos del haz libre; los segmentos indican los rangos al bloquearlo.}
\end{figure}

\subsection{Observaciones por condición}
\textbf{Lámpara.} Con el haz libre la diferencia se mantiene entre 3378 y 3385. Al interponer el objeto baja a 323--638. Es la condición con menor margen dentro de una misma prueba.

\textbf{Luz ambiental doméstica.} Con el haz libre se registraron valores \emph{Sin IR} de 4093--4095 y \emph{Con IR} de 171--177. Al interponer un objeto, el valor \emph{Con IR} pasó a 3491--3704 y la diferencia cayó a 391--604.

\textbf{Oscuridad.} La diferencia libre fue de 3913--3916 y la bloqueada de 319--739. En este montaje apagar las luces no impidió observar el haz IR.

\section{Umbral e histéresis propuestos}
Las lecturas están suficientemente separadas para plantear dos umbrales: uno para declarar la interrupción y otro para volver a declarar paso libre. De momento son valores \emph{propuestos}, no una detección automática ya validada.
\begin{center}
\begin{tabularx}{\linewidth}{@{}XlX@{}}
\toprule
Estado leído & Condición sobre $\Delta A$ & Acción\\
\midrule
Bloqueado & $\Delta A<1500$ & Marcar interrupción del haz\\
Zona intermedia & $1500\leq\Delta A\leq2500$ & Mantener el estado anterior\\
Libre & $\Delta A>2500$ & Marcar recuperación del haz\\
\bottomrule
\end{tabularx}
\end{center}

Según los resultados actuales, el mayor valor bloqueado (739) está 761 cuentas por debajo del umbral de 1500; el menor valor libre (3378) está 878 cuentas por encima del umbral de 2500. Son márgenes adecuados \emph{para los ensayos realizados}.

\begin{lstlisting}[language=C++,caption={Logica de estado que se propone probar a continuacion.}]
if (diferencia < 1500) {
    bloqueado = true;
} else if (diferencia > 2500) {
    bloqueado = false;
}
// Entre ambos umbrales se conserva el estado anterior.
\end{lstlisting}

\section{Limitaciones y siguientes pruebas}
El montaje responde claramente a un objeto situado entre el emisor y el receptor, pero todavía no puedo afirmar que detecte todos los pasos de una persona por una puerta. No he registrado distancia exacta, cantidad de repeticiones, velocidad de cruce ni tiempo de funcionamiento continuo.

Además, el firmware actual emplea dos esperas de 200 ms, diez muestras con 10 ms de pausa en cada estado y una espera final de 1000 ms. Por eso, sin contar pequeñas demoras de lectura y comunicación, el ciclo completo es de unos \SI{1.6}{\second}. Un paso rápido podría producirse entre medidas y no quedar registrado. Este será uno de los primeros puntos que tendré que mejorar.

El canal \emph{Sin IR} llega a 4095 en parte de los ensayos. Esto indica saturación de la escala del ADC y limita la información sobre intensidad luminosa en ese extremo, aunque no impidió distinguir los dos estados en las pruebas realizadas.

Las siguientes pruebas que quiero hacer son: variar la distancia y la alineación, introducir personas u objetos con diferentes velocidades, medir en luz solar directa, reducir el tiempo de adquisición y ensayar la histéresis contando interrupciones reales y falsas detecciones.

\section{Conclusión}
Después de corregir la polaridad del BPW40 y orientar el emisor y el receptor, he conseguido diferencias grandes y repetibles entre el haz libre y el bloqueado en las tres iluminaciones probadas. Esto valida la idea del sensor óptico en la protoboard y permite pasar al desarrollo de la detección automática. La instalación en la puerta, el aviso al móvil y la PCB final siguen pendientes.

\end{document}
